%%% LaTeX Template: Article/Thesis/etc. with colored headings and special fonts
%%%
%%% Source: http://www.howtotex.com/
% vim: set spell spelllang=es syntax=tex :

\documentclass[12pt]{article}
\usepackage[T1]{fontenc}
\usepackage{styles/apuntes-estilo}
\usepackage{styles/egyptian}
\usepackage{fancyhdr,lastpage}
\usepackage{hyperref}
\usepackage[inline]{enumitem}

\def\maketitle{

\makeatletter{
    \color{blue} \centering \huge \sc
    \textbf{
        Examen final libre 2022\\
        \small \vspace*{-8pt} \color{black}
        Parte práctica
        \vspace*{8pt}
    }\\
    \par
}
\makeatother

\makeatletter
% vim: set spell spelllang=es syntax=tex :
 {\centering \small 
    Introducción a la computación\\
    Departamento de Ingeniería de Computadoras \\
    Facultad de Informática - Universidad Nacional del Comahue \\
    \vspace{20pt} }
\makeatother

\vspace{-2.5cm}
\mbox{\hspace{-1cm}\includegraphics[width=3cm,height=3cm]{logos/uncoma.pdf}\hspace{12cm}
    \includegraphics[width=3cm,height=3cm]{logos/fai.pdf}}



}

% Custom headers and footers
\fancyhf{} % clear all header and footer fields
\fancypagestyle{plain}{\fancyhf{}}
\pagestyle{fancy}
\lhead{\footnotesize Primer examen parcial}
\rhead{\footnotesize \thepage\ }

\def\ti#1#2{\texttt{#1} & #2 \\ }

\begin{document}

\thispagestyle{empty}
\maketitle
\setlength{\parindent}{1pt}

\begin{enumerate}

    \item ¿Cuál es la representación de 42 (en base 7) en las siguientes bases?

        \begin{enumerate}

            \item Base 3.

            \item Base 2.

            \item Base 16.

        \end{enumerate}

    \item ¿Cuál es rango de representación del sistema de complemento a dos
        con 9 bytes?

    \item ¿Cual es la representación en \emph{punto fijo} y \emph{complemento
        a 2} de los siguientes números si se utilizan 4 bits para la parte
        entera y 4 para la parte fraccionaria? Indicar cual es el error de
        representación (la diferencia entre el valor original y el
        almacenado).

        \begin{enumerate}

            \item $7.4$

            \item $-7.4$

        \end{enumerate}

    \item Sea \textbf{0x418C0000} un número en base 16 (hexadecimal) que
        representa a un número real expresado en el sistema de \emph{Punto
        Flotante IEEE-754 de precisión simple (32 bits)}. Indique cuál es el
        valor en base 10 de dicho número real.

    \item ¿Que mensaje esta codificado en la siguiente secuencia de bytes
        expresados en hexadecimal: \textbf{46 69 6E 61 6C 20 64 65  20 49 43 2E}?
        \emph{Respetar mayúsculas y minúsculas.}

    \item Considerando la imagen que se muestra abajo, aplique un esquema de
        compresión que agrupa píxeles consecutivos de igual color y los
        reemplaza por una codificación ``\emph{cantidad/color}'', utilizando
        una codificación \textbf{4+2}, con cuatro bits para la cantidad y dos
        bits para el color. Asuma que 00 = blanco, 11 = negro y 01 = gris.

        \textit{Tenga en cuenta que al calcular la cantidad se debe considerar
        que las filas de la imagen son consecutivas. Es decir, si una fila
        termina con un grupo de pixeles un color y la siguiente comienza con
        pixeles del mismo color, se deben considerar todos como un solo grupo.
        \textbf{No olvide incluir la cabecera.}}

    \includegraphics[width=0.50\textwidth]{img/img_2020_1tpo-r.pdf}

    \item Si se tiene una imagen con 200 colores distintos ¿Cuál es la
        cantidad mínima de bits necesarios para la profundidad de color, de
        manera de poder representar la imagen sin pérdida de colores

    \item Dado el siguiente programa:

        \begin{verbatim}
     0: START:  LD  NUM
     1:         JZ  FIN
     2:         ST  OUT
     3:         SUB DOS
     4:         ST  NUM
     5:         JMP START
     6: FIN:    HLT
     7: NUM:    10
     8: DOS:    2
        \end{verbatim}

        \begin{enumerate}

            \item ¿Cuál es la salida del programa?


            \item ¿Cuál es el equivalente en código maquina binario de las
                instrucciones de las direcciones \textbf{0}, \textbf{1} y
                \textbf{5}?

        \end{enumerate}

    \item Dado el siguiente programa en binario:

        \begin{verbatim}
     0: 0100 0010
     1: 0111 1111
     2: 1010 0010
     3: 1111 1101
     4: 0010 0000
     5: 0000 0011
        \end{verbatim}

        ¿Cuál es la salida del programa?

    \item Si se desea extender la funcionalidad de la \emph{MCBE} para que
        tenga un total de 16 instrucciones distintas y la cantidad de celdas
        de memoria se extiende a 128 (incluyendo aquellas mapeadas a la
        \emph{entrada/salida} ¿Cuántos bits deberán tener el código de
        operación y el argumento?

\end{enumerate}

\subsection*{ \large\textbf{Anexo:} }

\textbf{Códigos de operación}

\begin{verbatim}
    010     LD
    011     ST
    100     ADD
    101     SUB
    110     JMP
    111     JZ
    001     HLT
    000     NOP
\end{verbatim}

\textbf{Tabla \emph{ASCII}:}

\begin{verbatim}
Dec Hex    Dec Hex    Dec Hex  Dec Hex  Dec Hex  Dec Hex   Dec Hex   Dec Hex
  0 00 NUL  16 10 DLE  32 20    48 30 0  64 40 @  80 50 P   96 60 `  112 70 p
  1 01 SOH  17 11 DC1  33 21 !  49 31 1  65 41 A  81 51 Q   97 61 a  113 71 q
  2 02 STX  18 12 DC2  34 22 "  50 32 2  66 42 B  82 52 R   98 62 b  114 72 r
  3 03 ETX  19 13 DC3  35 23 #  51 33 3  67 43 C  83 53 S   99 63 c  115 73 s
  4 04 EOT  20 14 DC4  36 24 $  52 34 4  68 44 D  84 54 T  100 64 d  116 74 t
  5 05 ENQ  21 15 NAK  37 25 %  53 35 5  69 45 E  85 55 U  101 65 e  117 75 u
  6 06 ACK  22 16 SYN  38 26 &  54 36 6  70 46 F  86 56 V  102 66 f  118 76 v
  7 07 BEL  23 17 ETB  39 27 '  55 37 7  71 47 G  87 57 W  103 67 g  119 77 w
  8 08 BS   24 18 CAN  40 28 (  56 38 8  72 48 H  88 58 X  104 68 h  120 78 x
  9 09 HT   25 19 EM   41 29 )  57 39 9  73 49 I  89 59 Y  105 69 i  121 79 y
 10 0A LF   26 1A SUB  42 2A *  58 3A :  74 4A J  90 5A Z  106 6A j  122 7A z
 11 0B VT   27 1B ESC  43 2B +  59 3B ;  75 4B K  91 5B [  107 6B k  123 7B {
 12 0C FF   28 1C FS   44 2C ,  60 3C <  76 4C L  92 5C \  108 6C l  124 7C |
 13 0D CR   29 1D GS   45 2D -  61 3D =  77 4D M  93 5D ]  109 6D m  125 7D }
 14 0E SO   30 1E RS   46 2E .  62 3E >  78 4E N  94 5E ^  110 6E n  126 7E ~
 15 0F SI   31 1F US   47 2F /  63 3F ?  79 4F O  95 5F _  111 6F o  127 7F DEL
\end{verbatim}

\end{document}
